\subsection{Formal realizability and extended Church's thesis}
\label{sec:p07-formal-realizability}

The external relation on numbers becomes an arithmetic interpretation
only after its computation language and its treatment of undefined
terms have been fixed. The following presentation specifies both.
The standard-numeral conclusion at the end is a metatheorem about
closed sentences, not an internal reflection principle for arithmetic.
For the wider account of realizability compare \cite{Troelstra1998};
the results here refer to the calculus and clauses displayed below.

\begin{hbdefinition}{def:p07-partial-arithmetic}{Arithmetic with partial application}
First expand HA definitionally by symbols for primitive-recursive
functions, with their defining equations. Use a fixed primitive-recursive
coding of finite computations, a computation predicate \(T(e,x,s)\),
and an output function \(U(s)\). Here \(T\) is represented by its
zero-valued primitive-recursive characteristic function. Choose the
ordinary program coding with effective parameter substitution; the
finite-code constructions used below are performed in this coding.

Adjoin a binary partial application term \(e\mathbin{\cdot}x\).
Write \(t\downarrow\) for \(t=t\), with equality read strictly:
both terms must denote. Variables denote, and functions and atomic
predicates are strict in their arguments. Quantifier instantiation
at \(t\) and existential introduction with witness \(t\) require
\(t\downarrow\). The other intuitionistic natural-deduction rules
retain their usual discharge and eigenvariable conditions. Equality
has variable reflexivity, symmetry, transitivity and the rule from
\(t=u\) and \(C(t)\) to \(C(u)\), with capture-avoiding
substitution. Axioms are universally closed; their instantiation at
partial terms retains the definedness guards. Equality does not assert
that a partial function is defined on equal arguments.
Primitive-recursive functions are total on defined arguments. The application axiom is
\[
 e\mathbin{\cdot}x=y\quad\longleftrightarrow\quad
 \exists s\bigl(T(e,x,s)\land U(s)=y\bigr).
\]
Together with HA's arithmetic axioms and induction for every formula
of this expanded language, these rules and axioms define
\(\mathsf{HA}^{\ast}\). There is no totality axiom for application.

An expression \(\lambda x.t\) in the proofs denotes the code of a
program evaluating \(t\), not a new object-language binder. Its
code is primitive recursive in the other free numerical parameters,
even when evaluating \(t\) may diverge. Application satisfies
\((\lambda x.t)\mathbin{\cdot}a\simeq t[a/x]\), where
\(u\simeq v\) means
\((u\downarrow\lor v\downarrow)\to u=v\).
Pairing and its projections \(\pi_0,\pi_1\) are the fixed total
primitive-recursive operations.
\end{hbdefinition}

\begin{hblemma}{thm:p07-partial-arithmetic-conservative}{proved=annex:p7.ect.1}{Eliminating partial terms}
\(\mathsf{HA}^{\ast}\) is conservative over the original HA for
sentences of the original arithmetic language. Finite program
composition, parameter substitution, and iteration have their
computation equations provable in \(\mathsf{HA}^{\ast}\).
\end{hblemma}

\begin{hbdefinition}{def:p07-formal-number-realizability}{Formal number realizability}
Define \(t\mathrel{\mathrm r}A\) for every term \(t\), including
partial terms, by recursion on the formula \(A\). Every clause
includes \(t\downarrow\). Beyond this conjunct, the clauses are:
\begin{enumerate}
\item an atom \(P\): \(P\) itself;
\item \(A\land B\):
      \(\pi_0t\mathrel{\mathrm r}A\land
        \pi_1t\mathrel{\mathrm r}B\);
\item \(A\lor B\): the tag \(\pi_0t\) is zero or one,
      and its corresponding payload \(\pi_1t\) realizes the
      selected formula;
\item \(A\to B\):
      \(\forall a(a\mathrel{\mathrm r}A\to
                      (t\mathbin{\cdot}a)\mathrel{\mathrm r}B)\);
\item \(\forall xA\):
      \(\forall x((t\mathbin{\cdot}x)\mathrel{\mathrm r}A)\);
\item \(\exists xA\):
      \(\pi_1t\mathrel{\mathrm r}A[\pi_0t/x]\).
\end{enumerate}
The disjunction clause is written without disjunction or existential
quantification in the interpreting formula. Let \(v(d)\) be zero
exactly for \(d=0,1\), and one otherwise. Its exact expansion is
\[
 v(\pi_0t)=0\ \land\
 (\pi_0t=0\to\pi_1t\mathrel{\mathrm r}A)\ \land\
 (\pi_0t=1\to\pi_1t\mathrel{\mathrm r}B).
\]
The arithmetic decidability of the tag proves its equivalence to the
selected-branch description. Falsity is the false atomic equation.
All substitutions are capture avoiding. In particular
\(t\mathrel{\mathrm r}A\to t\downarrow\) holds even when
\(A\) is an implication with an unrealizable antecedent.
\end{hbdefinition}

\begin{hbdefinition}{def:p07-almost-negative-ect}{Restricted extended Church's thesis}
Let \(\mathcal N\) be the formulas built from atoms and falsity
using conjunction, implication, and universal quantification in the
partial-term language. For \(D(x)\in\mathcal N\) and arbitrary
\(B(x,y)\), the partial-term scheme \(\mathsf{ECT}_0\) consists
of the universal closures of
\[
 \forall x\bigl(D(x)\to\exists yB(x,y)\bigr)
 \ \to\
 \exists e\,\forall x\bigl(D(x)\to
       (e\mathbin{\cdot}x\downarrow\land B(x,e\mathbin{\cdot}x))\bigr).
\]
Parameters are allowed; the variables bound in this display are fresh
where necessary. The restriction concerns the antecedent \(D\),
not the conclusion matrix \(B\).

In total primitive-recursive arithmetic, an \emph{almost-negative}
formula has no disjunction and has existential quantifiers only
immediately before atomic formulas. Each such existential atom is
equivalent in \(\mathsf{HA}^{\ast}\) to the definedness of the
program searching for its witness. Consequently its almost-negative
combinations have equivalents in \(\mathcal N\). Eliminating
partial application yields the familiar graph formulation
\[
 \forall x(D(x)\to\exists yB(x,y))\ \to\
 \exists e\,\forall x\bigl(D(x)\to
          \exists y(\chi(e,x,y)\land B(x,y))\bigr),
\]
where \(\chi\) is the displayed computation graph, not a new
primitive function. Conversely, eliminating a partial-term atom
produces a \(\Sigma^0_1\) formula, which can be written as an
existential atomic primitive-recursive equation. Thus the antecedent
classes correspond up to provable equivalence; they are not literally
the same strings in the two languages.
\index{extended Church thesis@extended Church's thesis}
\index{almost-negative formula}
\end{hbdefinition}

\begin{hblemma}{thm:p07-formal-realizability-lemmas}{proved=annex:p7.ect.2}{Canonical realizers and formal soundness}
Every formula \(t\mathrel{\mathrm r}A\) belongs to \(\mathcal N\).
For each \(D\in\mathcal N\) there is a term \(c_D\), with no
additional free parameters, such that \(\mathsf{HA}^{\ast}\) proves
\[
 (t\mathrel{\mathrm r}D\to D)
 \quad\hbox{and}\quad
 (D\to c_D\mathrel{\mathrm r}D).
\]
A derivation of \(A\) in \(\mathsf{HA}^{\ast}\) effectively yields
a term \(t\) and a derivation of \(t\mathrel{\mathrm r}A\) in
that same theory. Every instance of the restricted
\(\mathsf{ECT}_0\) also has a realizer whose realizing property is
provable in the base theory \(\mathsf{HA}^{\ast}\).
\end{hblemma}

\begin{hbtheorem}{thm:p07-ect-characterization}{proved=annex:p7.ect.3}{Formal number-realizability characterization}
For every formula \(A\), the explicitly specified theory satisfies
\[
 \mathsf{HA}^{\ast}+\mathsf{ECT}_0\ \vdash\
 A\longleftrightarrow\exists a(a\mathrel{\mathrm r}A).
\]
For every closed sentence \(A\) in its language,
\[
 \mathsf{HA}^{\ast}+\mathsf{ECT}_0\vdash A
 \quad\Longleftrightarrow\quad
 \mathsf{HA}^{\ast}\vdash\overline n\mathrel{\mathrm r}A.
\]
Here the right-hand side asserts the existence of such a proof for
some standard natural number \(n\). The numeral assertion is external. It does not assert a
primitive-recursive numeral-extraction bound or formalized
arithmetical reflection inside HA.
\end{hbtheorem}

\begin{hbtheorem}{thm:p07-unrestricted-ect-inconsistent}{proved=annex:p7.ect.4}{Unrestricted ECT}
Allowing arbitrary antecedents in the displayed
\(\mathsf{ECT}_0\) scheme makes \(\mathsf{HA}^{\ast}\)
inconsistent, even when its premise is restricted to unique
existence of the output. The contradiction requires no excluded-middle
axiom for arbitrary arithmetic formulas.
\end{hbtheorem}
